\chapter{Sức Khỏe Tâm Thần}
\markboth{Sức Khỏe Tâm Thần}{Mental Health}

\section{Không Phải Cứ Cười Là Ổn}
\begin{truyen}{Không Phải Cứ Cười Là Ổn}{Smiling Does Not Always Mean Fine}
\chuhoa{C}{ó những học sinh đi học đủ, nói chuyện bình thường, thậm chí còn pha trò.}
\textit{(Some students attend school, talk normally, and even joke around.)}

Nhưng bên trong các em đang kiệt quệ.
\textit{(Inside, however, they are exhausted.)}

Người lớn hay bỏ sót tín hiệu vì họ chỉ tìm dấu hiệu sụp đổ rõ ràng.
\textit{(Adults often miss the signs because they look only for dramatic collapse.)}

Trong khi nhiều đứa trẻ rất giỏi che giấu.
\textit{(Many children are very good at hiding.)}

Một học sinh có thể vẫn nộp bài đúng hạn và đồng thời mỗi tối đều nghĩ mình vô dụng.
\textit{(A student may still submit assignments on time while thinking every night that they are worthless.)}

Sức khỏe tâm thần không phải thứ chỉ được chú ý khi có khủng hoảng.
\textit{(Mental health is not something to notice only when crisis comes.)}

Nó cần được chăm trước khi vỡ.
\textit{(It must be cared for before it breaks.)}

\end{truyen}

\begin{baihoc}
Dạy học sinh và người lớn bỏ thói quen đánh giá ổn hay không chỉ bằng vẻ ngoài. Nhiều vết nứt nguy hiểm nhất nằm dưới gương mặt vẫn đang cười.
\textit{(Teach both students and adults to stop judging wellbeing by appearance alone. Many of the most dangerous cracks sit beneath a smiling face.)}
\end{baihoc}

\begin{ghinhoanh}
\textbf{Short English to remember:}

Teach both students and adults to stop judging wellbeing by appearance alone.

Many of the most dangerous cracks sit beneath a smiling face.

\end{ghinhoanh}

\begin{tuhocnhanh}
\textbf{Cau hoi tu hoc / Self-study questions}

1. Van de chinh la gi? \textit{What is the main problem?}

2. Cach xu ly tot hon la gi? \textit{What is a better action?}

3. Mot cau tieng Anh em muon nho la cau nao? \textit{Which English sentence do you want to remember?}
\end{tuhocnhanh}
\ngancach

\section{Kiệt Sức Không Phải Lười}
\begin{truyen}{Kiệt Sức Không Phải Lười}{Burnout Is Not Laziness}
\chuhoa{M}{ột học sinh từng chăm chỉ bỗng không còn muốn mở sách.}
\textit{(A once-hardworking student suddenly no longer wants to open a book.)}

Người lớn lập tức kết luận: lười.
\textit{(Adults immediately conclude: lazy.)}

Nhưng có khi em đã đi quá lâu trong áp lực mà không được nghỉ thật.
\textit{(But sometimes the student has been running under pressure too long without real rest.)}

Não và thân đã cạn.
\textit{(Mind and body are depleted.)}

Kiệt sức nhìn bên ngoài rất giống lười: chậm, nản, trì hoãn, né tránh.
\textit{(Burnout looks from the outside a lot like laziness: slow, discouraged, delaying, avoiding.)}

Nhưng gốc của nó khác.
\textit{(But the root is different.)}

Một bên là không muốn cố.
\textit{(One is not wanting to try.)}

Một bên là không còn lực để cố.
\textit{(The other is not having enough strength left to try.)}

Nếu nhầm hai thứ này, người lớn sẽ tiếp tục đánh vào chỗ đã nứt.
\textit{(If these two are confused, adults keep striking the place that is already cracked.)}

Và học sinh càng tệ hơn.
\textit{(Then the student gets worse.)}

\end{truyen}

\begin{baihoc}
Muốn giúp đúng, phải chẩn đoán đúng. Không phải đứa trẻ dừng lại nào cũng thiếu ý chí; có em chỉ đang cạn pin.
\textit{(To help correctly, you must diagnose correctly. Not every child who slows down lacks willpower; some are simply out of charge.)}
\end{baihoc}

\begin{ghinhoanh}
\textbf{Short English to remember:}

To help correctly, you must diagnose correctly.

Not every child who slows down lacks willpower; some are simply out of charge.

\end{ghinhoanh}

\begin{tuhocnhanh}
\textbf{Cau hoi tu hoc / Self-study questions}

1. Van de chinh la gi? \textit{What is the main problem?}

2. Cach xu ly tot hon la gi? \textit{What is a better action?}

3. Mot cau tieng Anh em muon nho la cau nao? \textit{Which English sentence do you want to remember?}
\end{tuhocnhanh}
\ngancach

\section{Xin Giúp Đỡ Không Làm Em Yếu Đi}
\begin{truyen}{Xin Giúp Đỡ Không Làm Em Yếu Đi}{Asking for Help Does Not Make You Weak}
\chuhoa{N}{hiều học sinh chịu đựng rất lâu vì sợ bị gắn mác yếu đuối, làm màu, hay có vấn đề.}
\textit{(Many students endure for a long time because they fear being labeled weak, dramatic, or problematic.)}

Các em chỉ đi tìm giúp đỡ khi mọi thứ đã quá nặng.
\textit{(They seek help only when everything has become too heavy.)}

Lúc đó việc hồi phục khó hơn nhiều.
\textit{(By then recovery is much harder.)}

Nếu đau bụng kéo dài, em đi khám.
\textit{(If stomach pain continues, you seek treatment.)}

Nếu tâm trí em tối dần, em cũng cần sự giúp đỡ chuyên môn.
\textit{(If your mind grows darker, you also need professional help.)}

Không có gì đáng xấu hổ ở đó.
\textit{(There is nothing shameful in that.)}

Một nền giáo dục tốt phải cho học sinh biết rằng cầu cứu đúng lúc là hành vi trưởng thành, không phải thất bại.
\textit{(A good educational culture must teach students that asking for help at the right time is maturity, not failure.)}

\end{truyen}

\begin{baihoc}
Hãy bình thường hóa việc tìm hỗ trợ tâm lý. Cái tôi im lặng nhiều khi nguy hiểm hơn chính vấn đề ban đầu.
\textit{(Normalize seeking mental-health support. A silent ego is often more dangerous than the initial problem itself.)}
\end{baihoc}

\begin{ghinhoanh}
\textbf{Short English to remember:}

Normalize seeking mental-health support.

A silent ego is often more dangerous than the initial problem itself.

\end{ghinhoanh}

\begin{tuhocnhanh}
\textbf{Cau hoi tu hoc / Self-study questions}

1. Van de chinh la gi? \textit{What is the main problem?}

2. Cach xu ly tot hon la gi? \textit{What is a better action?}

3. Mot cau tieng Anh em muon nho la cau nao? \textit{Which English sentence do you want to remember?}
\end{tuhocnhanh}
\ngancach

\section{Nghỉ Ngơi Không Phải Tội Lỗi}
\begin{truyen}{Nghỉ Ngơi Không Phải Tội Lỗi}{Rest Is Not a Sin}
\chuhoa{C}{ó những học sinh cứ nghỉ là thấy cắn rứt.}
\textit{(Some students feel guilty whenever they rest.)}

Không học một tối cũng thấy mình tụt lại.
\textit{(Missing one evening of study already feels like falling behind.)}

Đó không còn là chăm chỉ nữa.
\textit{(That is no longer diligence.)}

Đó là mối quan hệ bệnh lý với năng suất.
\textit{(It is an unhealthy relationship with productivity.)}

Não người không vận hành như máy.
\textit{(The human mind does not operate like a machine.)}

Nó cần khoảng dừng để hợp nhất kiến thức, để hồi phục cảm xúc, để không cháy rụi.
\textit{(It needs pauses to consolidate knowledge, recover emotionally, and avoid burning out.)}

Học sinh cần được dạy rằng nghỉ đúng cách không cản trở thành công.
\textit{(Students need to be taught that proper rest does not obstruct success.)}

Nó là điều kiện để thành công không giết mình.
\textit{(It is what keeps success from destroying you.)}

\end{truyen}

\begin{baihoc}
Nghỉ ngơi có kỷ luật là một phần của học tập có kỷ luật. Đừng dạy trẻ em tự hào vì kiệt sức.
\textit{(Disciplined rest is part of disciplined study. Do not teach children to feel proud of exhaustion.)}
\end{baihoc}

\begin{ghinhoanh}
\textbf{Short English to remember:}

Disciplined rest is part of disciplined study.

Do not teach children to feel proud of exhaustion.

\end{ghinhoanh}

\begin{tuhocnhanh}
\textbf{Cau hoi tu hoc / Self-study questions}

1. Van de chinh la gi? \textit{What is the main problem?}

2. Cach xu ly tot hon la gi? \textit{What is a better action?}

3. Mot cau tieng Anh em muon nho la cau nao? \textit{Which English sentence do you want to remember?}
\end{tuhocnhanh}
\ngancach

\section{Người Học Trò Vẫn Đi Học Đều Nhưng Đã Rỗng Dần}
\begin{truyen}{Người Học Trò Vẫn Đi Học Đều Nhưng Đã Rỗng Dần}{The Student Who Still Attended but Was Slowly Emptying Out}
\chuhoa{C}{ó học sinh vẫn đến lớp đầy đủ, làm bài đủ, cười khi cần, nhưng bên trong đã rỗng đi từng ngày.}
\textit{(Some students still attend class, finish assignments, and smile when needed, while inside they are emptying out day by day.)}

Người lớn dễ bị đánh lừa bởi sự đúng giờ và bề ngoài gọn gàng.
\textit{(Adults are easily misled by punctuality and neat appearance.)}

Nhưng một tâm trí kiệt quệ vẫn có thể vận hành thêm một thời gian bằng quán tính.
\textit{(Yet an exhausted mind can keep functioning for a while through inertia.)}

Nếu chỉ chờ đến lúc một học sinh gục hẳn mới gọi đó là vấn đề, ta đã chậm rất lâu rồi.
\textit{(If we wait until a student completely collapses before calling it a problem, we are already very late.)}

\end{truyen}

\begin{baihoc}
Sức khỏe tâm thần cần được quan sát ở giai đoạn nứt, không phải đợi đến giai đoạn vỡ.
\textit{(Mental health must be noticed at the cracking stage, not only at the breaking stage.)}
\end{baihoc}

\begin{ghinhoanh}
\textbf{Short English to remember:}

Mental health must be noticed at the cracking stage, not only at the breaking stage.

\end{ghinhoanh}

\begin{tuhocnhanh}
\textbf{Cau hoi tu hoc / Self-study questions}

1. Van de chinh la gi? \textit{What is the main problem?}

2. Cach xu ly tot hon la gi? \textit{What is a better action?}

3. Mot cau tieng Anh em muon nho la cau nao? \textit{Which English sentence do you want to remember?}
\end{tuhocnhanh}
\ngancach

\section{Ông Đồ Già Và Cơn U Uất Mùa Thi}
\begin{truyen}{Ông Đồ Già Và Cơn U Uất Mùa Thi}{The Old Tutor and the Exam-Season Melancholy}
\chuhoa{N}{gày xưa mỗi mùa khoa cử đến, không ít sĩ tử mất ăn mất ngủ, dễ cáu bẳn, chán nản, và sợ hãi quá mức.}
\textit{(In the old examination seasons, many candidates lost sleep and appetite, became irritable, despondent, and excessively fearful.)}

Người đời khi đó chưa gọi tên là kiệt quệ tâm lý.
\textit{(People then did not name it psychological burnout.)}

Nhưng triệu chứng thì vẫn là những gì học sinh hôm nay đang trải qua.
\textit{(Yet the symptoms were much like what students experience today.)}

Một ông đồ già thường bắt học trò nghỉ hẳn một buổi chiều trước kỳ thi để đi bộ quanh đồng và nói chuyện trời đất.
\textit{(An old tutor would make his students take one full afternoon off before the exam to walk through the fields and talk about ordinary things.)}

Ông bảo: "Đầu óc căng quá thì chữ nghĩa dính vào nhau thành bùn.
\textit{(He said, "When the mind is too tight, words stick together like mud.)}

Phải để nó có chỗ thở."
\textit{(It must be given space to breathe.")}

\end{truyen}

\begin{baihoc}
Vấn đề tâm lý trong học tập không phải chuyện mới của thời hiện đại. Nó chỉ có tên gọi mới hơn mà thôi.
\textit{(Psychological struggle in study is not a new modern problem. It simply has newer names now.)}
\end{baihoc}

\begin{ghinhoanh}
\textbf{Short English to remember:}

Psychological struggle in study is not a new modern problem.

It simply has newer names now.

\end{ghinhoanh}

\begin{tuhocnhanh}
\textbf{Cau hoi tu hoc / Self-study questions}

1. Van de chinh la gi? \textit{What is the main problem?}

2. Cach xu ly tot hon la gi? \textit{What is a better action?}

3. Mot cau tieng Anh em muon nho la cau nao? \textit{Which English sentence do you want to remember?}
\end{tuhocnhanh}
\ngancach

\section{Cái Mệt Không Nói Thành Lời}
\begin{truyen}{Cái Mệt Không Nói Thành Lời}{The Fatigue Without Words}
\chuhoa{N}{hiều học sinh không biết mô tả mình đang bị gì.}
\textit{(Many students do not know how to describe what is happening inside them.)}

Các em chỉ nói: ``Em mệt.''
\textit{(They simply say, ``I am tired.'')}

Mệt có thể là thiếu ngủ.
\textit{(Tired may mean lack of sleep.)}

Cũng có thể là lo âu, là quá tải, là mất hứng sống, là áp lực kéo dài.
\textit{(It may also mean anxiety, overload, loss of joy, or prolonged pressure.)}

Khi từ vựng cảm xúc quá nghèo, học sinh rất khó cầu cứu đúng thứ mình cần.
\textit{(When emotional vocabulary is too poor, students struggle to ask for the help they actually need.)}

Dạy một đứa trẻ gọi tên được trạng thái của mình là đang cho nó thêm một công cụ sống sót.
\textit{(Teaching a child to name their inner state gives them one more survival tool.)}

\end{truyen}

\begin{baihoc}
Giáo dục cảm xúc không phải thứ phụ. Nó giúp học sinh phân biệt mệt, buồn, kiệt sức, hoảng loạn, và biết khi nào cần nói ra.
\textit{(Emotional education is not an extra. It helps students distinguish tiredness, sadness, burnout, panic, and know when to speak.)}
\end{baihoc}

\begin{ghinhoanh}
\textbf{Short English to remember:}

Emotional education is not an extra.

It helps students distinguish tiredness, sadness, burnout, panic, and know when to speak.

\end{ghinhoanh}

\begin{tuhocnhanh}
\textbf{Cau hoi tu hoc / Self-study questions}

1. Van de chinh la gi? \textit{What is the main problem?}

2. Cach xu ly tot hon la gi? \textit{What is a better action?}

3. Mot cau tieng Anh em muon nho la cau nao? \textit{Which English sentence do you want to remember?}
\end{tuhocnhanh}
\ngancach

\section{Người Đứng Đầu Lớp Cũng Có Thể Muốn Biến Mất}
\begin{truyen}{Người Đứng Đầu Lớp Cũng Có Thể Muốn Biến Mất}{The Top Student Can Also Want to Disappear}
\chuhoa{T}{hành tích tốt không miễn trừ ai khỏi đau tinh thần.}
\textit{(Good performance does not immunize anyone against mental pain.)}

Một học sinh đứng đầu lớp vẫn có thể về nhà và cảm thấy trống rỗng, vô nghĩa, hoặc muốn biến mất khỏi mọi kỳ vọng.
\textit{(A top student can still go home and feel empty, meaningless, or wish to disappear from every expectation.)}

Chính vì người đó học tốt nên xung quanh càng dễ chủ quan: nó vẫn ổn mà.
\textit{(Precisely because the student performs well, people around are more likely to be careless: he seems fine.)}

Đó là một sai lầm đắt.
\textit{(That is a costly mistake.)}

Năng lực không phải áo giáp trước tổn thương tâm lý.
\textit{(Competence is not armor against psychological harm.)}

\end{truyen}

\begin{baihoc}
Đừng đánh giá sức khỏe tâm thần bằng kết quả học tập. Một bảng điểm đẹp có thể che rất nhiều nứt gãy.
\textit{(Do not judge mental health by academic results. A beautiful report card can hide many fractures.)}
\end{baihoc}

\begin{ghinhoanh}
\textbf{Short English to remember:}

Do not judge mental health by academic results.

A beautiful report card can hide many fractures.

\end{ghinhoanh}

\begin{tuhocnhanh}
\textbf{Cau hoi tu hoc / Self-study questions}

1. Van de chinh la gi? \textit{What is the main problem?}

2. Cach xu ly tot hon la gi? \textit{What is a better action?}

3. Mot cau tieng Anh em muon nho la cau nao? \textit{Which English sentence do you want to remember?}
\end{tuhocnhanh}
\ngancach

\section{Khi Cơ Thể Kêu Thay Tâm Trí}
\begin{truyen}{Khi Cơ Thể Kêu Thay Tâm Trí}{When the Body Speaks for the Mind}
\chuhoa{C}{ó học sinh đau bụng trước giờ kiểm tra, tim đập mạnh, run tay, toát mồ hôi, nhưng không hề nghĩ đó là chuyện tâm lý.}
\textit{(Some students get stomach pain before tests, racing hearts, shaking hands, and sweating, yet do not realize it may be psychological.)}

Các em tưởng mình yếu người.
\textit{(They think their body is weak.)}

Thực ra có khi cơ thể đang nói thay cho một đầu óc quá căng.
\textit{(In fact, the body may be speaking on behalf of an overloaded mind.)}

Nếu chỉ chữa phần ngọn mà không nhìn phần gốc, nỗi khổ sẽ đổi hình chứ không biến mất.
\textit{(If only the surface is treated while the root is ignored, the suffering changes shape but does not disappear.)}

Học sinh cần biết tâm trí và cơ thể không tách nhau sạch như trong sách vở.
\textit{(Students need to know that mind and body are not separated as neatly as textbooks suggest.)}

\end{truyen}

\begin{baihoc}
Một số tín hiệu cơ thể là lời nhắc rằng tâm trí đang quá tải. Biết đọc những dấu hiệu đó là kỹ năng tự bảo vệ.
\textit{(Some bodily symptoms are reminders that the mind is overloaded. Learning to read them is a self-protection skill.)}
\end{baihoc}

\begin{ghinhoanh}
\textbf{Short English to remember:}

Some bodily symptoms are reminders that the mind is overloaded.

Learning to read them is a self-protection skill.

\end{ghinhoanh}

\begin{tuhocnhanh}
\textbf{Cau hoi tu hoc / Self-study questions}

1. Van de chinh la gi? \textit{What is the main problem?}

2. Cach xu ly tot hon la gi? \textit{What is a better action?}

3. Mot cau tieng Anh em muon nho la cau nao? \textit{Which English sentence do you want to remember?}
\end{tuhocnhanh}
\ngancach

\section{Người Lớn Lắng Nghe Một Lần Đúng Cách}
\begin{truyen}{Người Lớn Lắng Nghe Một Lần Đúng Cách}{One Adult Who Listens Properly}
\chuhoa{C}{ó học sinh đã nhiều lần định nói mình không ổn, nhưng mỗi lần mở lời đều bị chặn bằng câu: ``Nghĩ ít thôi là hết.''}
\textit{(A student tried many times to say they were not okay, but every attempt was cut off with, ``Just think less and it will pass.'')}

Đến khi gặp một cô giáo chỉ hỏi: ``Em muốn cô ngồi nghe hay muốn cô tìm người giúp?'' em mới bật khóc.
\textit{(Only when a teacher asked, ``Do you want me to listen, or do you want me to help find support?'' did the student finally break down in tears.)}

Được lắng nghe tử tế không phải phép màu.
\textit{(Being listened to properly is not magic.)}

Nhưng nó thường là cánh cửa đầu tiên để một người chịu bước ra khỏi bóng tối.
\textit{(But it is often the first door through which someone steps out of darkness.)}

Không phải người lớn nào cũng cần giải quyết ngay.
\textit{(Not every adult needs to solve things immediately.)}

Nhiều khi việc đầu tiên là đừng làm học sinh thấy mình bị coi nhẹ.
\textit{(Often the first task is simply not to make the student feel minimized.)}

\end{truyen}

\begin{baihoc}
Trong hỗ trợ tâm lý, cách nghe quan trọng không kém cách nói. Một câu đúng lúc có thể giữ người ở lại với đời.
\textit{(In mental-health support, the way you listen matters as much as the way you speak. One timely sentence can help keep someone in life.)}
\end{baihoc}

\begin{ghinhoanh}
\textbf{Short English to remember:}

In mental-health support, the way you listen matters as much as the way you speak.

One timely sentence can help keep someone in life.

\end{ghinhoanh}

\begin{tuhocnhanh}
\textbf{Cau hoi tu hoc / Self-study questions}

1. Van de chinh la gi? \textit{What is the main problem?}

2. Cach xu ly tot hon la gi? \textit{What is a better action?}

3. Mot cau tieng Anh em muon nho la cau nao? \textit{Which English sentence do you want to remember?}
\end{tuhocnhanh}
